\documentclass[12pt,a4paper,oneside]{report}

% ============================================================
% ECOTWIN PROJECT REPORT
% ============================================================

% -------------------- Page layout ----------------------------
\usepackage[
    a4paper,
    left=1.25in,
    right=1.0in,
    top=1.0in,
    bottom=1.0in
]{geometry}

% -------------------- Fonts / encoding -----------------------
\usepackage[T1]{fontenc}
\usepackage{lmodern}
\usepackage{microtype}

% -------------------- Mathematics ----------------------------
\usepackage{amsmath}
\usepackage{amssymb}
\usepackage{amsfonts}

% -------------------- Figures --------------------------------
\usepackage{graphicx}
\usepackage{float}
\usepackage{subcaption}
\usepackage{wrapfig}

% -------------------- Tables ---------------------------------
\usepackage{booktabs}
\usepackage{longtable}
\usepackage{array}
\usepackage{multirow}
\usepackage{tabularx}
\usepackage{makecell}

% -------------------- Lists / units ---------------------------
\usepackage{enumitem}
\usepackage{siunitx}

% -------------------- Colors ---------------------------------
\usepackage{xcolor}

% -------------------- Code -----------------------------------
\usepackage{listings}

% -------------------- Diagrams -------------------------------
\usepackage{tikz}
\usetikzlibrary{
    arrows.meta,
    positioning,
    shapes.geometric,
    fit,
    backgrounds
}

% -------------------- Headers / links ------------------------
\usepackage{fancyhdr}
\usepackage[
    colorlinks=true,
    linkcolor=black,
    citecolor=black,
    urlcolor=blue
]{hyperref}
\usepackage[nameinlink,noabbrev]{cleveref}


% ============================================================
% DOCUMENT INFORMATION
% ============================================================

\newcommand{\projecttitle}{EcoTwin}
\newcommand{\projectsubtitle}{
Ecosystem Analog Search and Forecasting Using Geospatial Foundation Models
}
\newcommand{\studentname}{Your Name}
\newcommand{\studentroll}{Your Roll Number}
\newcommand{\supervisorname}{Supervisor Name}
\newcommand{\departmentname}{Computer Science and Engineering}
\newcommand{\institutionname}{International Institute of Information Technology Hyderabad}
\newcommand{\reportyear}{2026}

% ============================================================
% CUSTOM COMMANDS
% ============================================================

\newcommand{\code}[1]{\texttt{#1}}
\newcommand{\topk}{Top-K}

\newcommand{\placeholderfigure}[3][0.8\textwidth]{
    \begin{figure}[H]
        \centering
        \fbox{
            \begin{minipage}[c][#2][c]{#1}
                \centering
                \textit{Figure placeholder}\\[5pt]
                #3
            \end{minipage}
        }
        \caption{#3}
    \end{figure}
}

% ============================================================
% CODE STYLE
% ============================================================

\definecolor{codegray}{gray}{0.95}
\definecolor{commentgreen}{RGB}{30,110,30}

\lstset{
    backgroundcolor=\color{codegray},
    basicstyle=\ttfamily\small,
    breaklines=true,
    frame=single,
    numbers=left,
    numberstyle=\tiny,
    keywordstyle=\bfseries,
    commentstyle=\color{commentgreen},
    showstringspaces=false,
    tabsize=4
}

% ============================================================
% HEADER / FOOTER
% ============================================================

\pagestyle{fancy}
\fancyhf{}
\fancyhead[L]{\projecttitle}
\fancyhead[R]{\leftmark}
\fancyfoot[C]{\thepage}
\renewcommand{\headrulewidth}{0.4pt}
\renewcommand{\footrulewidth}{0pt}
\setlength{\headheight}{18pt}

\setlength{\parindent}{0pt}
\setlength{\parskip}{6pt}
\emergencystretch=3em
\renewcommand{\arraystretch}{1.2}

% ============================================================
% DOCUMENT
% ============================================================


\begin{document}

% ============================================================
% TITLE PAGE
% ============================================================

\begin{titlepage}
\begin{center}
\vspace*{1cm}

{\Large \textbf{\institutionname}}

\vspace{0.5cm}
{\large \departmentname}

\vspace{2cm}

{\Huge \textbf{\projecttitle}}

\vspace{0.7cm}

{\Large \textbf{\projectsubtitle}}

\vspace{2cm}

{\large Project Report}

\vspace{1.5cm}

Submitted by

\vspace{0.4cm}
{\Large \textbf{\studentname}}

\vspace{0.2cm}
{\large \studentroll}

\vspace{1.5cm}

Under the guidance of

\vspace{0.4cm}
{\Large \textbf{\supervisorname}}

\vfill
{\large \reportyear}

\end{center}
\end{titlepage}

% % ============================================================
% % CERTIFICATE
% % ============================================================

% \chapter*{Certificate}
% \addcontentsline{toc}{chapter}{Certificate}

% This is to certify that the project report entitled
% \textbf{\projectsubtitle} is a record of the work carried out by
% \textbf{\studentname} under my supervision and guidance.

% \vspace{2cm}

% \begin{flushright}
% \textbf{\supervisorname}\\
% Supervisor
% \end{flushright}

% \clearpage

% % ============================================================
% % DECLARATION
% % ============================================================

% \chapter*{Declaration}
% \addcontentsline{toc}{chapter}{Declaration}

% I hereby declare that this project report entitled
% \textbf{\projectsubtitle} is the result of the work carried out
% as part of the project requirements.

% All external sources, datasets, models, and software components
% used in the project have been appropriately acknowledged.

% \vspace{2cm}

% \begin{flushright}
% \textbf{\studentname}\\
% \studentroll
% \end{flushright}

% \clearpage

% % ============================================================
% % ACKNOWLEDGEMENTS
% % ============================================================

% \chapter*{Acknowledgements}
% \addcontentsline{toc}{chapter}{Acknowledgements}

% We would like to express our sincere gratitude to our supervisor,
% faculty members, colleagues, and everyone who contributed to the
% development of EcoTwin.

% We also acknowledge the providers of the HydroLAKES, Sentinel-2,
% Dynamic World, Google Earth Engine, and Prithvi resources, as well
% as the open-source PostgreSQL, PostGIS, pgvector, Python, FastAPI,
% React, Leaflet, Redis, Celery, and Docker ecosystems used during
% development.

% \clearpage

% ============================================================
% ABSTRACT
% ============================================================

\chapter*{Abstract}
\addcontentsline{toc}{chapter}{Abstract}

EcoTwin is an end-to-end geospatial ecosystem analog discovery
platform designed around the question: \emph{which ecosystems
elsewhere exhibit similar spatial structure, and what can their
historical trajectories tell us about a target ecosystem?}

The system combines a static lake inventory with multi-year
Sentinel-2 observations, spatial tiling, environmental indicators,
geospatial foundation-model embeddings, vector similarity search,
and temporal ecosystem profiles. HydroLAKES polygons provide the
initial ecosystem boundaries and metadata. For each selected lake,
the pipeline projects and buffers the geometry, constructs a
manageable imagery grid, retrieves annual Sentinel-2 composites
from Google Earth Engine, stores the resulting GeoTIFFs locally,
and subsequently subdivides the merged imagery into approximately
one-kilometre analysis cells. Cells with insufficient lake coverage
are discarded.

The retained image cells are converted into the six-band input
required by the Prithvi geospatial foundation model. The resulting
768-dimensional representations are stored in PostgreSQL using
pgvector. A query ecosystem can then be compared against the
indexed collection using cosine similarity, producing a ranked set
of analogous ecosystems. The web application exposes these results
through an interactive map and profile views, together with
historical NDVI, NDWI, and NBR trajectories.

The project was implemented iteratively over approximately two
months. A substantial part of the engineering work involved making
large-scale Earth Engine extraction reliable, handling historical
Sentinel-2 collection differences, reducing export sizes, merging
local raster tiles, controlling concurrency, validating embeddings,
and designing a database capable of preserving both spatial
geometries and vector features. The forecasting layer is designed
as an analog-based, non-parametric extension in which historical
trajectories of retrieved twins provide empirical evidence for
future ecosystem behaviour rather than requiring synthetic image
generation.

\textbf{Keywords:}
Geospatial Foundation Models, Prithvi, Remote Sensing, Sentinel-2,
HydroLAKES, Ecosystem Analogs, Vector Search, pgvector, NDVI, NDWI,
NBR, Temporal Analysis

\clearpage

% ============================================================
% CONTENTS
% ============================================================

\tableofcontents
\clearpage
\listoffigures
\clearpage
\listoftables
\clearpage

% ============================================================
% CHAPTER 1
% ============================================================

\chapter{Introduction}
\label{chap:introduction}

\section{Background}

Earth observation satellites provide repeated measurements of
ecosystems at spatial scales that cannot be obtained economically
through field observations alone. Sentinel-2 is particularly useful
for this purpose because it provides multispectral optical
observations suitable for vegetation, water, land-cover, and
surface-condition analysis.

A conventional ecosystem-analysis workflow normally begins with a
specific question such as vegetation change, water extent, or land
cover. EcoTwin approaches the problem from a complementary
direction: instead of analysing a lake only in isolation, it first
asks whether another ecosystem in a large reference collection
looks similar in its spatial-spectral representation.

The project therefore treats ecosystem comparison as a retrieval
problem. Satellite observations are converted into learned feature
vectors, and those vectors are indexed so that an ecosystem can be
compared against a large collection without repeatedly processing
the raw imagery.

\section{Problem Statement}

The practical problem is to identify geographically distributed
ecosystems that are similar to a selected target ecosystem using
information contained in satellite observations.

The problem has several engineering constraints:

\begin{enumerate}
    \item lake boundaries are irregular polygons rather than
    fixed-size image squares;
    \item large satellite scenes are expensive to download and store;
    \item historical Sentinel-2 collections do not expose identical
    quality bands across all years;
    \item cloud and invalid-pixel contamination must be controlled;
    \item large ecosystems may require multiple image tiles;
    \item the embedding pipeline must process many spatial regions;
    \item embeddings must remain associated with their source
    geometries and lake/year metadata; and
    \item online similarity search must be substantially cheaper
    than recomputing embeddings.
\end{enumerate}

\section{Core EcoTwin Hypothesis}

EcoTwin is based on the operational hypothesis:

\begin{quote}
\emph{Ecosystems elsewhere that are structurally similar in
geospatial feature space can provide useful analogs for understanding
the spatial and historical behaviour of a target ecosystem.}
\end{quote}

The hypothesis separates two tasks:

\begin{enumerate}
    \item \textbf{Analog discovery:} find ecosystems that resemble
    the target in the learned spatial-spectral representation.
    \item \textbf{Trajectory analysis:} compare historical
    environmental indicators of the target and its analogs.
\end{enumerate}

\section{Research Objectives}

The project objectives are:

\begin{enumerate}
    \item construct a scalable geospatial data-ingestion pipeline;
    \item represent lake environments using standardized satellite
    image regions;
    \item generate fixed-dimensional Prithvi embeddings;
    \item store and index embeddings using PostgreSQL and pgvector;
    \item retrieve ranked analog ecosystems using cosine similarity;
    \item expose lake boundaries and analogs through an interactive
    geospatial interface;
    \item compare historical NDVI, NDWI, and NBR profiles; and
    \item extend analog retrieval into an empirical forecasting
    mechanism based on historical twin trajectories.
\end{enumerate}

\section{Project Evolution During Development}

The implementation evolved considerably from the initial conceptual
design. The early design considered standardized ecosystem blocks
of approximately $5\times5$ km. During implementation, the data
engineering problem became more important than the initial geometric
assumption. Large lake polygons and Earth Engine export limits made
it preferable to separate \emph{imagery acquisition tiles} from
\emph{analysis cells}.

The current implementation therefore uses a lake-centred buffered
imagery extraction stage followed by local raster merging and
approximately $1\times1$ km analysis cells. The larger imagery tiles
are an I/O and Earth Engine management mechanism; the smaller cells
are the units from which embeddings are generated. This distinction
allows large lakes to be processed without forcing the foundation
model to operate directly on irregular polygons.

\section{Contributions}

The project contributions are:

\begin{itemize}
    \item an end-to-end lake-to-embedding-to-analog pipeline;
    \item a robust satellite-image acquisition strategy using
    Google Earth Engine and local raster processing;
    \item polygon-aware spatial subdivision and coverage filtering;
    \item Prithvi-based six-band geospatial representations;
    \item persistent 768-dimensional vector storage;
    \item cosine-similarity analog retrieval;
    \item temporal indicator comparison;
    \item a FastAPI backend and React/Leaflet frontend; and
    \item an engineering workflow for validation and database
    consolidation across multiple development environments.
\end{itemize}

\section{Report Organization}

Chapter~\ref{chap:data} describes the source datasets and data
engineering process. Chapter~\ref{chap:methodology} presents the
complete scientific pipeline. Chapter~\ref{chap:implementation}
describes the software and database implementation.
Chapter~\ref{chap:results} documents the current experimental
outcomes and engineering observations. Chapter~\ref{chap:discussion}
discusses interpretation and limitations. Chapter~\ref{chap:forecast}
describes the analog-based forecasting layer. Chapter~\ref{chap:conclusion}
summarizes the project and future work.

% ============================================================
% CHAPTER 2
% ============================================================

\chapter{Data Sources and Data Engineering}
\label{chap:data}

\section{Overview}

EcoTwin combines static geographic information with repeated
satellite observations. The static layer defines \emph{where} an
ecosystem is, while the satellite layer defines \emph{what the
ecosystem looks like through time}.

The data-engineering pipeline was designed so that expensive
satellite processing is performed offline. The online application
then operates primarily on database records, precomputed features,
and indexed embeddings.

\section{HydroLAKES Ecosystem Inventory}

HydroLAKES provides lake polygons and associated metadata. The
database representation used by EcoTwin retains fields including:

\begin{itemize}
    \item \code{hylak\_id} / lake identifier;
    \item lake name;
    \item country;
    \item lake area;
    \item elevation;
    \item pour-point latitude and longitude;
    \item polygon geometry;
    \item lake type;
    \item average depth;
    \item total volume; and
    \item watershed area.
\end{itemize}

The complete HydroLAKES-derived inventory is much larger than the
cohort used for the computationally expensive processing stages.
The development pipeline was narrowed to a selected processing
cohort of lakes so that multi-year imagery and embeddings could be
generated within available storage, Earth Engine, and compute
budgets.

\begin{table}[H]
\centering
\caption{HydroLAKES-derived information used by EcoTwin.}
\label{tab:hydrolakes_fields}
\begin{tabularx}{\textwidth}{l l X}
\toprule
Field & Type & Role in EcoTwin \\
\midrule
Lake ID & Integer & Primary ecosystem identifier \\
Lake name & Text & Human-readable identification \\
Country & Text & Geographic filtering \\
Lake area & Numeric & Size and metadata analysis \\
Elevation & Numeric & Environmental context \\
Pour latitude/longitude & Numeric & Geographic reference \\
Geometry & PostGIS geometry & Spatial processing and mapping \\
Lake type & Categorical & Ecosystem metadata \\
Average depth & Numeric & Descriptive context \\
Total volume & Numeric & Descriptive context \\
Watershed area & Numeric & Catchment context \\
\bottomrule
\end{tabularx}
\end{table}

\section{Sentinel-2}

Sentinel-2 multispectral imagery is the primary remote-sensing
input. Google Earth Engine is used for collection filtering,
compositing, masking, and export. The principal collection used
during development was:

\begin{center}
\code{COPERNICUS/S2\_SR\_HARMONIZED}
\end{center}

For years or locations where the surface-reflectance harmonized
collection did not provide observations, the pipeline introduced
a fallback to:

\begin{center}
\code{COPERNICUS/S2}
\end{center}

This fallback was particularly important for early Sentinel-2 years.

\section{Temporal Coverage}

The initial analytical design focused on a 2018--2025 operational
window. During implementation, ingestion was extended to earlier
Sentinel-2 observations where available, resulting in an attempted
2016--2025 extraction window.

The practical coverage is therefore lake- and year-dependent.
Early years contain fewer observations and, in some cases, require
the Level-1C fallback collection. The pipeline records the
collection actually used so that missing or fallback observations
are not silently treated as equivalent to the primary surface
reflectance product.

\section{Satellite Bands}

The current embedding input uses six bands:

\begin{center}
\code{B2, B3, B4, B5, B6, B7}
\end{center}

These correspond to the blue, green, red, narrow near-infrared, and
two short-wave infrared channels used by the selected Prithvi
configuration.

The bands are selected consistently from the exported GeoTIFFs
before model inference.

\section{Dynamic World}

Dynamic World was evaluated as a complementary land-cover/context
source. It can provide 10 m land-use/land-cover information and can
be used to augment interpretation of the spatial analogs.

In the current EcoTwin architecture, Sentinel-2 is the principal
source for the embedding image, while Dynamic World is treated as a
contextual layer rather than as a replacement for the Prithvi input.

\section{Environmental Indices}

EcoTwin maintains temporal environmental indicators for ecosystem
comparison.

\subsection{NDVI}

\begin{equation}
NDVI =
\frac{NIR-Red}{NIR+Red}
\end{equation}

NDVI provides a vegetation-sensitive signal and is useful for
studying seasonal and interannual vegetation behaviour.

\subsection{NDWI}

The water-sensitive index is computed as:

\begin{equation}
NDWI =
\frac{Green-NIR}{Green+NIR}
\end{equation}

It provides a complementary signal for surface-water behaviour.

\subsection{NBR}

The Normalized Burn Ratio is:

\begin{equation}
NBR =
\frac{NIR-SWIR}{NIR+SWIR}
\end{equation}

NBR provides information related to surface disturbance and
burn-related spectral change.

The exact band mapping used for each index is kept consistent with
the final feature-extraction implementation.

\section{Data Quality Requirements}

A valid ecosystem record requires:

\begin{enumerate}
    \item a valid lake polygon;
    \item a valid year;
    \item a successful imagery query;
    \item a usable composite;
    \item valid exported raster data;
    \item valid spatial cells after coverage filtering; and
    \item a successfully generated embedding for each retained cell.
\end{enumerate}

The pipeline explicitly distinguishes successful, failed, skipped,
and already-completed work so that long-running jobs can be resumed.

% ============================================================
% CHAPTER 3
% ============================================================

\chapter{Detailed EcoTwin Methodology}
\label{chap:methodology}

\section{End-to-End Pipeline}

The implemented workflow is:

\begin{equation}
\begin{aligned}
Lake\ Inventory
&\rightarrow Boundary\ Validation\\
&\rightarrow Buffering
\rightarrow Imagery\ Tiling\\
&\rightarrow Sentinel\text{-}2\ Composite
\rightarrow GeoTIFF\\
&\rightarrow Local\ Mosaic
\rightarrow 1\,km\ Cells\\
&\rightarrow Coverage\ Filter
\rightarrow Prithvi\\
&\rightarrow 768D\ Embeddings
\rightarrow pgvector\\
&\rightarrow Cosine\ Search
\rightarrow Analogs\\
&\rightarrow Temporal\ Profiles
\rightarrow Forecast
\end{aligned}
\end{equation}

\section{Step 1: Lake Boundary Retrieval}

The pipeline begins by reading the target lake polygon from
PostGIS. Keeping the original geometry in a spatial database avoids
repeated conversion between shapefiles and application-specific
formats.

The geometry is checked for validity and transformed into an
appropriate projected CRS before distance-based operations are
performed.

\section{Step 2: Buffering}

The lake polygon is expanded using a configurable buffer. During
development, a \SI{1200}{\meter} buffer was used for the imagery
extraction stage.

The buffer is important because the ecosystem representation is not
restricted to the water polygon itself. The surrounding land,
shoreline, vegetation, and local landscape context contribute to
the spatial-spectral representation.

\section{Step 3: Projected Grid Generation}

The buffered polygon is projected into a metric CRS. A regular
imagery grid is then generated around the buffered geometry.

The imagery grid is deliberately different from the final
analysis grid. During the latest extraction implementation, an
approximately \SI{8500}{\meter} tile size was used to balance Earth
Engine request size, memory use, and local processing cost.

A lake may therefore produce one, two, or several imagery tiles
depending on its size and shape.

\section{Step 4: Earth Engine Collection Selection}

For each lake-year combination, the pipeline first queries the
preferred surface-reflectance collection. If no usable observation
exists, the system falls back to the Level-1C collection.

This decision is made before the export request and is recorded in
the processing metadata.

\begin{table}[H]
\centering
\caption{Sentinel-2 collection strategy.}
\label{tab:collection_strategy}
\begin{tabularx}{\textwidth}{c l X}
\toprule
Priority & Collection & Purpose \\
\midrule
1 & \code{COPERNICUS/S2\_SR\_HARMONIZED} &
Preferred surface reflectance \\
2 & \code{COPERNICUS/S2} &
Fallback for early or unavailable surface-reflectance years \\
\bottomrule
\end{tabularx}
\end{table}

\section{Step 5: Cloud and Invalid-Pixel Handling}

Cloud contamination is one of the main sources of noise in optical
remote sensing. The current preprocessing pipeline uses
cloud-probability information and class-based masking where the
required bands are available.

The development configuration used a cloud-probability threshold
of approximately $40\%$. Pixels classified as unusable are masked
before compositing.

The implementation had to account for historical collection
differences. Early Sentinel-2 products exposed different quality
bands; code that assumed the presence of \code{QA60},
\code{MSK\_CLASSI\_OPAQUE}, or \code{MSK\_CLASSI\_CIRRUS} could fail
when the selected collection did not expose the expected band.
The final pipeline therefore treats cloud masking as a
collection-aware operation rather than assuming one universal
quality-band schema.

\section{Step 6: Annual Composite Construction}

For each lake-year-region, the filtered observations are combined
into an annual composite.

The median operation is used to reduce the influence of residual
clouds, outliers, and individual acquisition anomalies:

\begin{equation}
I_{annual}(p,b)
=
median\left(
I_1(p,b),I_2(p,b),\ldots,I_n(p,b)
\right)
\end{equation}

where $p$ is a pixel and $b$ is a selected spectral band.

The composite is clipped to the requested geographic extent before
export.

\section{Step 7: GeoTIFF Export}

The selected six bands are exported as local GeoTIFF data. The
development pipeline used 10 m spatial output for the selected
Sentinel-2 bands and stored the raster in floating-point form.

Local storage has two major advantages:

\begin{enumerate}
    \item the same Earth Engine composite does not need to be
    recomputed for every downstream analysis cell; and
    \item embedding generation can be repeated without repeatedly
    consuming Earth Engine export capacity.
\end{enumerate}

\section{Step 8: Tile Completion and Resumability}

Each imagery tile is represented by metadata including lake ID,
year, tile index, geometry/bounding box, status, and local image
path.

The pipeline checks whether a tile has already completed before
starting a new Earth Engine request. This makes the large batch
pipeline resumable after failures or interruptions.

\section{Step 9: Local Raster Merging}

After all tiles belonging to a lake-year are available, the local
processing stage merges them into a single raster mosaic.

Rasterio is used for reading, reprojection/metadata handling,
mosaicking, and writing. Compression settings are applied to
reduce storage requirements while preserving the numerical data
needed for model inference.

The merged image becomes the stable source from which smaller
analysis cells are extracted.

\section{Step 10: Analysis-Cell Generation}

The merged lake image is intersected with a regular approximately
$1\times1$ km analysis grid.

For each cell, the lake coverage is computed:

\begin{equation}
Coverage(c)
=
\frac{Area(c\cap Lake)}
{Area(c)}
\times 100
\end{equation}

Only cells above the configured minimum coverage threshold are
retained. The development configuration used a threshold of
approximately $15\%$.

This prevents cells that only touch a small corner of a lake from
being treated as representative ecosystem regions.

\section{Step 11: Model Input Preparation}

Each retained analysis cell is read from the local merged GeoTIFF.
The six selected bands are converted into the model's expected
tensor layout, numerical range, and spatial dimensions.

The processing stage therefore separates:

\begin{enumerate}
    \item geographic geometry;
    \item raster extraction;
    \item model preprocessing; and
    \item model inference.
\end{enumerate}

This separation makes it possible to change the model without
re-downloading the underlying satellite data.

\section{Step 12: Prithvi Embedding Generation}

Prithvi is used as the geospatial foundation model. Each retained
analysis cell is transformed into a fixed-length feature vector.

For an input image $X$:

\begin{equation}
f(X) \rightarrow \mathbf{z},
\qquad
\mathbf{z}\in\mathbb{R}^{768}
\end{equation}

The six-band configuration used by EcoTwin is based on
\code{B2--B7}, corresponding to the selected blue, green, red,
narrow-NIR, SWIR1, and SWIR2 channels.

The inference pipeline supports batched processing. During the
development runs, a batch size of eight was used in CPU-based
inference.

\section{Step 13: Embedding Validation}

Before insertion into the database, every embedding is checked for:

\begin{itemize}
    \item expected dimensionality;
    \item missing or null values;
    \item numerical validity;
    \item correct lake association;
    \item correct year association; and
    \item correct spatial-cell association.
\end{itemize}

A failed inference is not inserted as a valid vector.

\section{Step 14: Vector Persistence}

The final embedding is stored in PostgreSQL using pgvector.

The conceptual relationship is:

\begin{equation}
Lake
\rightarrow
Year
\rightarrow
Spatial\ Cell
\rightarrow
Embedding
\end{equation}

This preserves the provenance required to explain why an ecosystem
was retrieved as an analog.

\section{Step 15: Similarity Search}

For a query embedding $\mathbf{q}$ and candidate embedding
$\mathbf{x}$, EcoTwin uses cosine similarity:

\begin{equation}
sim(\mathbf{q},\mathbf{x})
=
\frac{\mathbf{q}^{T}\mathbf{x}}
{\|\mathbf{q}\|_2\|\mathbf{x}\|_2}
\end{equation}

The candidates are ranked by decreasing similarity:

\begin{equation}
s_{(1)}\ge s_{(2)}\ge\cdots\ge s_{(K)}
\end{equation}

The highest-ranked ecosystems become the analog set displayed by
the application.

\section{Cell-to-Lake Interpretation}

Because a lake can contain multiple analysis cells, the retrieval
layer must preserve the relationship between cell-level embeddings
and lake-level identity.

This design has two advantages:

\begin{enumerate}
    \item it preserves local spatial detail instead of collapsing a
    large lake into one opaque vector too early; and
    \item it allows retrieved similarities to be traced back to the
    actual spatial region that produced the embedding.
\end{enumerate}

Where a single lake-level representation is required, cell
representations can be aggregated using a coverage- or area-aware
strategy. The cell-level vectors remain the underlying evidence.

\section{Step 16: Temporal Profile Generation}

For the query lake and retrieved analogs, historical environmental
records are assembled into time series.

The primary indicators are:

\begin{itemize}
    \item NDVI;
    \item NDWI; and
    \item NBR.
\end{itemize}

The interface can show the indicators independently or as a
multi-index comparison.

\section{Step 17: Analog-Based Forecasting Concept}

EcoTwin's forecasting design is not based on generating synthetic
future satellite images. Instead, it uses the observed historical
trajectory of an analogous ecosystem as empirical evidence.

Let the target indicator be $y_t$ and an analog indicator be
$a_t$. A simple trajectory-offset formulation is:

\begin{equation}
\Delta_t = y_t-a_t
\end{equation}

The future analog trajectory can then be transferred to the target
through an estimated historical offset:

\begin{equation}
\hat{y}_{t+h}
=
a_{t+h}
+
\widehat{\Delta}
\end{equation}

For multiple analogs, a weighted combination can be used:

\begin{equation}
\hat{y}_{t+h}
=
\sum_{i=1}^{K}
w_i
\left(
a_{i,t+h}+\Delta_i
\right)
\end{equation}

with weights related to the similarity scores and normalized so that

\begin{equation}
\sum_{i=1}^{K}w_i=1.
\end{equation}

This formulation keeps the forecast grounded in historically
observed ecosystem trajectories.

\section{Why the Analog Architecture is Explainable}

A parametric time-series model produces a prediction through learned
parameters that can be difficult to connect to a particular
ecosystem.

EcoTwin instead provides an explicit chain:

\begin{equation}
\text{Target}
\rightarrow
\text{Retrieved Twin}
\rightarrow
\text{Similarity}
\rightarrow
\text{Historical Trajectory}
\rightarrow
\text{Projected Trajectory}
\end{equation}

This makes it possible to inspect the analogs responsible for a
forecast and compare their historical behaviour with the target.

% ============================================================
% CHAPTER 4
% ============================================================

\chapter{System Architecture and Implementation}
\label{chap:implementation}

\section{Architecture Overview}

EcoTwin is implemented as a multi-component system. The offline
pipeline performs the computationally expensive work, while the
online application retrieves precomputed data.

\placeholderfigure[0.95\textwidth]{10cm}{
Recommended figure: EcoTwin end-to-end architecture showing
HydroLAKES, Google Earth Engine, local raster processing, Prithvi,
PostgreSQL/PostGIS/pgvector, FastAPI, and React/Leaflet.
}

\section{Technology Stack}

\begin{table}[H]
\centering
\caption{EcoTwin implementation stack.}
\label{tab:stack}
\begin{tabularx}{\textwidth}{l l X}
\toprule
Layer & Technology & Responsibility \\
\midrule
Frontend & React, Vite, TypeScript & Interactive application \\
Mapping & Leaflet / React-Leaflet & Spatial visualization \\
Backend & FastAPI, Python & REST services and orchestration \\
Database & PostgreSQL & Relational persistence \\
Spatial database & PostGIS & Polygon and spatial queries \\
Vector search & pgvector & Embedding storage and similarity search \\
Remote sensing & Google Earth Engine & Satellite filtering/compositing \\
Raster processing & Rasterio / Python & Local GeoTIFF processing \\
ML inference & Prithvi & Geospatial embeddings \\
Task queue & Celery / Redis & Asynchronous processing \\
Monitoring & Flower & Celery task monitoring \\
Administration & pgAdmin & Database administration \\
Containerization & Docker Compose & Reproducible local deployment \\
\bottomrule
\end{tabularx}
\end{table}

\section{Offline and Online Separation}

The system is divided into two logical paths.

\subsection{Offline Path}

The offline path performs:

\begin{enumerate}
    \item lake preparation;
    \item imagery acquisition;
    \item raster merging;
    \item spatial-cell extraction;
    \item embedding generation;
    \item temporal-feature preparation; and
    \item vector indexing.
\end{enumerate}

\subsection{Online Path}

The online path performs:

\begin{enumerate}
    \item lake lookup;
    \item map rendering;
    \item similarity query;
    \item analog profile retrieval;
    \item temporal-series retrieval; and
    \item report generation.
\end{enumerate}

This architecture avoids running Earth Engine or the foundation
model during every user request.

\section{PostgreSQL and PostGIS}

PostgreSQL is the central persistent store. PostGIS stores the
original lake geometry, analysis-cell geometry, and spatial metadata.

Spatial information remains in the database rather than being
converted into plain latitude/longitude tables. This allows
intersection, containment, area, and map-query operations to be
performed directly in the database.

\section{pgvector}

The pgvector extension stores the 768-dimensional Prithvi vectors.

The vector table maintains the association between:

\begin{equation}
Embedding
\leftrightarrow
Subregion
\leftrightarrow
Lake
\leftrightarrow
Year
\end{equation}

Cosine-distance/similarity operations can therefore be executed
directly against the indexed database rather than exporting vectors
to CSV or recomputing features.

\section{Core Database Entities}

The implementation evolved around entities including:

\begin{itemize}
    \item \code{lakes};
    \item \code{lake\_tiles};
    \item \code{lake\_images};
    \item \code{sub\_regions};
    \item \code{sub\_region\_features}; and
    \item embedding/vector records.
\end{itemize}

\subsection{Lakes}

Stores the master ecosystem identity, metadata, and PostGIS
geometry.

\subsection{Lake Images}

Stores the merged or processed imagery associated with a lake-year,
including local path and processing status.

\subsection{Lake Tiles}

Stores the intermediate imagery tiles generated for a lake-year,
including tile index, geometry, bounding box, status, and image path.

\subsection{Sub-Regions}

Stores the smaller analysis cells generated from the merged lake
image, including geometry and coverage information.

\subsection{Sub-Region Features}

Stores model-derived features and embeddings associated with
individual spatial cells.

\section{Backend Services}

The FastAPI backend provides services for:

\begin{itemize}
    \item lake metadata lookup;
    \item nearest lake/grid-cell lookup;
    \item analog similarity search;
    \item temporal feature retrieval;
    \item forecast retrieval;
    \item report generation; and
    \item health/system status.
\end{itemize}

The API separates data retrieval from heavy batch processing. Long
operations are handled asynchronously where appropriate.

\section{Frontend}

The frontend is implemented using React, Vite, and TypeScript.
The map uses Leaflet/React-Leaflet.

The main user interaction is:

\begin{enumerate}
    \item select or search for a lake;
    \item display its geographic boundary;
    \item inspect lake metadata;
    \item request analogous ecosystems;
    \item visualize the analog locations;
    \item inspect similarity scores;
    \item inspect analog profiles; and
    \item compare temporal indicators and forecast information.
\end{enumerate}

\section{Map Rendering and Geometry Provenance}

An important implementation requirement is that displayed lake
boundaries must come from the same PostGIS geometry used for the
analysis. Analog markers are derived from their stored geographic
coordinates and lake identities.

This prevents a common failure mode in which the analytical result
and the displayed map use different coordinate transformations or
different numbering systems.

\section{Asynchronous Processing}

Celery and Redis are used for background processing. This is
important because imagery downloads, raster merging, and model
inference can take substantially longer than a normal HTTP request.

Flower is used to inspect worker activity and task progress during
development.

\section{Dockerized Development Environment}

The local development environment was organized as a multi-service
Docker Compose deployment containing services such as:

\begin{itemize}
    \item PostgreSQL;
    \item Redis;
    \item FastAPI backend;
    \item Celery worker;
    \item pipeline worker;
    \item Flower; and
    \item pgAdmin.
\end{itemize}

The containerized design makes service dependencies explicit and
allows the same architecture to be reproduced on another machine.

\section{Configuration Management}

The processing pipeline uses environment/configuration values for
items such as:

\begin{itemize}
    \item Earth Engine project/authentication;
    \item date ranges;
    \item cloud thresholds;
    \item imagery buffer distance;
    \item imagery tile size;
    \item analysis-cell size;
    \item coverage threshold;
    \item local storage paths;
    \item database credentials; and
    \item worker/batch settings.
\end{itemize}

This prevents experiment-specific values from being hard-coded
throughout the processing code.

% ============================================================
% CHAPTER 5
% ============================================================

\chapter{Experimental Results and Engineering Evaluation}
\label{chap:results}

\section{Current Processing Cohort}

During development, a processing cohort of approximately 829 lakes
was targeted for multi-year processing. The wider HydroLAKES-derived
Indian inventory is substantially larger, so the 829-lake cohort
represents the computationally processed subset rather than the
entire source inventory.

\begin{table}[H]
\centering
\caption{Current EcoTwin processing configuration.}
\label{tab:current_config}
\begin{tabular}{ll}
\toprule
Parameter & Current configuration \\
\midrule
Lake source & HydroLAKES \\
Target processing cohort & $\sim$829 lakes \\
Imagery window & 2016--2025 attempted \\
Primary Sentinel collection & S2 SR Harmonized \\
Fallback collection & S2 Level-1C \\
Imagery buffer & 1200 m \\
Imagery tile size & approximately 8500 m \\
Analysis cell & approximately 1 km $\times$ 1 km \\
Coverage threshold & approximately 15\% \\
Embedding model & Prithvi \\
Embedding dimension & 768 \\
Input bands & B2--B7 \\
Similarity metric & Cosine similarity \\
\bottomrule
\end{tabular}
\end{table}

\section{Temporal Availability Observations}

The Earth Engine collection checks demonstrated that early years
have substantially less usable surface-reflectance coverage than
recent years. In particular, the 2016--2017 period required special
handling, while later years provided much denser observation sets.

This is an important distinction between the \emph{requested}
temporal window and the \emph{available valid observations}. EcoTwin
therefore stores year-level processing status rather than assuming
that every lake has a complete ten-year sequence.

\section{Embedding Pipeline Performance}

The embedding stage was implemented as batch inference. During
development, inference logs showed CPU execution with a batch size
of eight.

The pipeline is designed so that:

\begin{equation}
N_{embeddings}
=
N_{valid\ cells}
\end{equation}

rather than

\begin{equation}
N_{embeddings}
=
N_{lakes}.
\end{equation}

This distinction is fundamental to EcoTwin: a large lake can
generate several spatial embeddings, and the vector index therefore
contains spatially localized ecosystem representations.

\section{Similarity Search Demonstration}

A representative similarity-analysis run used a large candidate
collection and returned a ranked set of analogs. In a previously
observed application run, approximately 12,458 candidate lake/region
records were compared, the interface returned the top 10 analogs,
and the best displayed similarity was approximately 0.942.

These values should be treated as \emph{current demonstration
results} rather than final benchmark claims until the final
production database is frozen and the experiment is repeated under
a documented configuration.

\begin{table}[H]
\centering
\caption{Representative similarity-search demonstration.}
\label{tab:similarity_demo}
\begin{tabular}{ll}
\toprule
Metric & Observed value \\
\midrule
Candidate records & $\sim$12,458 \\
Returned analogs & 10 \\
Best displayed similarity & $\sim$0.942 \\
Similarity metric & Cosine similarity \\
\bottomrule
\end{tabular}
\end{table}

\section{Interpretation of Similarity Scores}

A similarity score indicates closeness in the learned representation
space. It does not directly imply identical ecological processes.

A high score can therefore be interpreted as:

\begin{quote}
The target and analog contain similar spatial-spectral information
under the selected satellite bands, preprocessing procedure, and
Prithvi representation.
\end{quote}

Ecological interpretation requires additional evidence from lake
metadata, spatial maps, and temporal environmental indicators.

\section{Temporal Profile Results}

The temporal layer provides a second axis of validation. Two lakes
can be spatially similar but exhibit different historical
trajectories, while two lakes with different geographic locations
may show similar seasonal behaviour.

EcoTwin therefore presents:

\begin{itemize}
    \item NDVI trajectories;
    \item NDWI trajectories;
    \item NBR trajectories;
    \item observation availability;
    \item yearly summaries; and
    \item comparison against retrieved analogs.
\end{itemize}

\section{Processing Reliability}

The system was intentionally designed to tolerate partial failure.

A lake-year can be in states such as:

\begin{itemize}
    \item pending;
    \item processing;
    \item completed;
    \item failed; or
    \item skipped because output already exists.
\end{itemize}

This makes large multi-year extraction jobs restartable instead of
requiring a complete restart after a single failed Earth Engine
request.

\section{Engineering Problems Encountered}

\subsection{Earth Engine Memory Limits}

Large requests sometimes produced Earth Engine memory-limit errors.
The solution was to reduce the size of each export region and to
move part of the processing to local raster operations.

\subsection{Earth Engine Concurrency Limits}

Parallel workers caused HTTP 429 errors when too many requests were
issued simultaneously. Connection-pool saturation was also observed.

The pipeline was consequently changed toward smaller controlled
worker counts, resumable jobs, and local reuse of already-generated
imagery.

\subsection{Historical Band Differences}

The early Sentinel-2 Level-1C fallback exposed a different set of
quality bands from assumptions made for other collections. For
example, requests that selected \code{MSK\_CLASSI\_CIRRUS},
\code{MSK\_CLASSI\_OPAQUE}, or \code{QA60} could fail depending on
the actual image schema.

This led to collection-aware preprocessing rather than a single
hard-coded cloud-mask implementation.

\subsection{Large Raster Storage}

Merged rasters can become much larger than individual source tiles.
Compression, floating-point type control, tiled GeoTIFF writing,
and local reuse were therefore important.

\subsection{Disk Constraints}

During development, the local root filesystem reached 100\%
utilization, causing raster and Docker operations to fail with
``No space left on device''. This demonstrated that the data
pipeline must account for intermediate storage, not only final
dataset size.

\subsection{Network and Authentication Constraints}

Earth Engine authentication and access depend on external network
services. Authentication failures and inability to reach
OAuth endpoints temporarily prevented pipeline execution.

The implementation therefore separates authentication from
per-lake processing and records failures instead of silently
producing incomplete datasets.

\section{Data Migration Between Team Members}

Because the project was developed by multiple team members, a
database consolidation procedure was designed.

Each developer can create a PostgreSQL custom-format dump. The dump
is restored into a temporary database, validated through a Python
ETL process, and then merged into the production database.

The procedure preserves:

\begin{itemize}
    \item PostGIS geometries;
    \item pgvector embeddings;
    \item indexes;
    \item constraints; and
    \item PostgreSQL data types.
\end{itemize}

This approach is safer than CSV-based migration because CSV does not
naturally preserve PostGIS geometry and vector types.

% ============================================================
% CHAPTER 6
% ============================================================

\chapter{Discussion}
\label{chap:discussion}

\section{Why Foundation-Model Embeddings}

Traditional similarity systems require manually specifying which
variables should define similarity. EcoTwin instead uses a learned
geospatial representation derived from multispectral imagery.

This changes the problem from:

\[
Similarity =
f(\text{manually selected attributes})
\]

to:

\[
Similarity =
f(\text{learned spatial-spectral representation})
\]

The advantage is that spatial context can be represented jointly
rather than through a small set of manually selected scalar
attributes.

\section{Why Spatial Cells Instead of One Lake Image}

A lake is not spatially homogeneous. Shoreline vegetation,
surrounding land cover, built-up regions, agricultural land,
mountain slopes, and water surfaces can coexist in the same
analysis area.

Using smaller cells preserves this heterogeneity.

The cell-level representation also improves explainability because
a retrieved similarity can be associated with a geographic region
rather than only with an entire lake.

\section{Why Local Raster Processing}

Downloading the complete India-scale satellite archive would be
impractical for the project. EcoTwin therefore uses an
``extract once, reuse many times'' strategy:

\[
Earth\ Engine
\rightarrow
Local\ GeoTIFF
\rightarrow
Multiple\ Downstream\ Features
\]

This substantially reduces repeated remote requests.

\section{Why Vector Database Search}

Without vector persistence, every query would require model
inference on the candidate ecosystem collection.

With pgvector:

\[
Query\ Image
\rightarrow
One\ Embedding
\rightarrow
Vector\ Search
\rightarrow
Top\text{-}K
\]

The expensive foundation-model inference is therefore moved into
the offline indexing stage.

\section{Spatial Similarity versus Temporal Similarity}

EcoTwin deliberately keeps spatial and temporal similarity
conceptually separate.

Spatial similarity answers:

\begin{quote}
Which ecosystems look similar to the target in the learned
representation?
\end{quote}

Temporal similarity answers:

\begin{quote}
Do those ecosystems exhibit similar historical behaviour?
\end{quote}

A robust analog should ideally provide useful evidence on both axes.

\section{Strengths}

The implemented system provides:

\begin{itemize}
    \item scalable offline preprocessing;
    \item resumable Earth Engine extraction;
    \item explicit spatial provenance;
    \item foundation-model representations;
    \item persistent vector search;
    \item temporal contextualization;
    \item interactive geospatial visualization; and
    \item a path toward explainable analog-based forecasting.
\end{itemize}

\section{Limitations}

\subsection{Observation Availability}

The ten-year target window does not guarantee ten complete years
for every lake. Historical collection availability and cloud
filtering can create gaps.

\subsection{Optical Cloud Dependence}

Sentinel-2 is an optical sensor. Cloud cover and atmospheric
conditions can reduce valid observations.

\subsection{Embedding Semantics}

A high embedding similarity does not guarantee ecological
equivalence. The embedding is a representation of the selected
satellite inputs and model.

\subsection{Spatial Grid Dependence}

Changing the analysis-cell size or coverage threshold changes the
resulting embeddings. The spatial discretization is therefore a
model parameter.

\subsection{Model Input Dependence}

Changing bands, preprocessing, normalization, temporal compositing,
or image extent can change the embedding space.

\subsection{Forecasting Dependence on Analogs}

Analog forecasting is only as strong as the retrieved historical
twins. Poor analogs lead to poor forecasts, even if the mathematical
forecasting procedure itself is correct.

\section{Validation Strategy}

The system should be validated at four levels:

\begin{enumerate}
    \item \textbf{Data validation:} check image completeness,
    geometry validity, and temporal availability.
    \item \textbf{Representation validation:} verify embedding
    dimensionality and numerical validity.
    \item \textbf{Retrieval validation:} inspect similarity rankings
    and compare retrieved maps and metadata.
    \item \textbf{Temporal validation:} determine whether high-ranked
    analogs also provide useful historical trajectories.
\end{enumerate}

Independent ecological or field observations would provide the
strongest external validation but were not available as a complete
ground-truth source during the present development cycle.

% ============================================================
% CHAPTER 7
% ============================================================

\chapter{Analog-Based Forecasting}
\label{chap:forecast}

\section{Forecasting Philosophy}

The EcoTwin forecasting layer follows the same analog architecture
as the retrieval system. Instead of treating forecasting as an
independent machine-learning problem, it reuses the evidence already
identified by similarity search.

The logic is:

\begin{equation}
\text{Current Target}
\rightarrow
\text{Historical Twin}
\rightarrow
\text{Observed Future Trajectory}
\rightarrow
\text{Target Forecast}
\end{equation}

\section{Historical Twin Alignment}

Suppose the target ecosystem has an indicator sequence
$Y_t$ and an analog has sequence $A_t$. A historical overlap period
is used to estimate how the target differs from the analog.

A simple offset estimate is:

\begin{equation}
\widehat{\Delta}
=
\frac{1}{T}
\sum_{t=1}^{T}
(Y_t-A_t)
\end{equation}

The future analog trajectory is then adjusted using this offset.

\section{Multiple Analog Forecast}

For $K$ analogs, each analog receives a weight $w_i$ derived from
its similarity. A generic normalization is:

\begin{equation}
w_i =
\frac{s_i}
{\sum_{j=1}^{K}s_j}
\end{equation}

for non-negative similarity scores $s_i$.

The predicted trajectory becomes:

\begin{equation}
\widehat{Y}_{t+h}
=
\sum_{i=1}^{K}
w_i
\left(
A_{i,t+h}+\widehat{\Delta}_i
\right).
\end{equation}

The application can therefore explain a prediction through the
analog lakes that contributed to it.

\section{Forecasted Variables}

The principal targets are:

\begin{itemize}
    \item NDVI;
    \item NDWI;
    \item NBR; and
    \item a combined ecosystem indicator derived from these series,
    where implemented.
\end{itemize}

\section{Forecast Uncertainty}

The spread among analog predictions can provide an empirical
uncertainty signal.

For example, if the analog trajectories are widely separated,
forecast confidence should be lower than when the trajectories
closely agree.

This is not equivalent to a formally calibrated probabilistic
forecast and should be described as an analog-dispersion measure
unless a formal uncertainty model is added.

\section{Forecast Evaluation}

Once a sufficiently long historical sequence is available, a
hindcasting experiment can be performed:

\begin{enumerate}
    \item hide the latest observed period;
    \item retrieve analogs using only information available before
    the cutoff;
    \item generate a forecast;
    \item compare the forecast against the held-out observations.
\end{enumerate}

Recommended metrics include:

\begin{equation}
MAE =
\frac{1}{n}\sum_{i=1}^{n}|y_i-\hat{y}_i|
\end{equation}

and

\begin{equation}
RMSE =
\sqrt{
\frac{1}{n}
\sum_{i=1}^{n}
(y_i-\hat{y}_i)^2
}.
\end{equation}

The evaluation should be performed separately for NDVI, NDWI, and
NBR rather than reporting only a single aggregate number.

% ============================================================
% CHAPTER 8
% ============================================================

\chapter{Conclusion and Future Work}
\label{chap:conclusion}

\section{Conclusion}

EcoTwin evolved from a conceptual ecosystem-comparison idea into
an end-to-end geospatial retrieval system.

The final architecture separates the problem into several stages:
a lake inventory defines the ecosystem boundary; Google Earth Engine
provides multi-year Sentinel-2 observations; a buffered and tiled
raster pipeline makes large exports manageable; local mosaicking
creates reusable lake-year images; one-kilometre analysis cells
provide spatially consistent units; Prithvi converts these cells into
768-dimensional representations; pgvector makes those
representations searchable; and the web application exposes the
retrieved analogs together with spatial and temporal context.

The most significant engineering contribution of the project is
therefore not a single model component but the integration of
geospatial data engineering, foundation-model inference, vector
retrieval, temporal analysis, and interactive visualization into one
reproducible workflow.

The development experience also demonstrated that large-scale
Earth-observation systems require careful handling of external
service limits, historical data differences, raster storage,
concurrency, resumability, and database provenance.

\section{Future Work}

Future work should focus on:

\begin{enumerate}
    \item completing and systematically evaluating the analog-based
    forecasting layer;
    \item performing temporal hindcasting experiments;
    \item calibrating forecast uncertainty using analog dispersion
    and historical error;
    \item adding independent ecological validation;
    \item improving the lake-level aggregation strategy;
    \item experimenting with alternative cell sizes and coverage
    thresholds;
    \item incorporating additional environmental variables;
    \item improving temporal compositing;
    \item supporting larger ecosystem inventories;
    \item optimizing vector indexing for larger collections;
    \item improving asynchronous pipeline orchestration; and
    \item deploying a shared production database and hosted
    processing infrastructure.
\end{enumerate}

\section{Final Remarks}

EcoTwin provides a practical bridge between modern geospatial
foundation models and ecosystem-scale comparative analysis. Its
central idea is to use historically observed ecosystems as
computationally retrievable analogs rather than treating every
ecosystem as an isolated forecasting problem.

% ============================================================
% REFERENCES
% ============================================================

\begin{thebibliography}{99}

\bibitem{gorelick2017gee}
N. Gorelick, M. Hancher, M. Dixon, S. Ilyushchenko, D. Thau, and
R. Moore, ``Google Earth Engine: Planetary-scale geospatial analysis
for everyone,'' \textit{Remote Sensing of Environment}, vol. 202,
pp. 18--27, 2017.

\bibitem{sentinel2}
European Space Agency, ``Sentinel-2 Mission,'' Copernicus
documentation.

\bibitem{hydrolakes}
HydroLAKES, Global lake database and associated lake-boundary
information used as the static ecosystem inventory.

\bibitem{prithvi}
NASA and IBM, ``Prithvi-EO: Geospatial foundation model for Earth
observation,'' official model documentation and associated
publication.

\bibitem{dynamicworld}
C. F. Brown \textit{et al.}, ``Dynamic World, near real-time global
10 m land use land cover mapping,'' 2022.

\bibitem{pettorelli2005ndvi}
N. Pettorelli, J. O. Vik, A. Mysterud, J.-M. Gaillard, C. J. Tucker,
and N. C. Stenseth, ``Using the satellite-derived NDVI to assess
ecological responses to environmental change,'' \textit{Trends in
Ecology \& Evolution}, vol. 20, no. 9, pp. 503--510, 2005.

\bibitem{huete2002modis}
A. Huete, K. Didan, T. Miura, E. P. Rodriguez, X. Gao, and
L. G. Ferreira, ``Overview of the radiometric and biophysical
performance of the MODIS vegetation indices,'' \textit{Remote
Sensing of Environment}, vol. 83, no. 1--2, pp. 195--213, 2002.

\end{thebibliography}

% ============================================================
% APPENDIX A
% ============================================================

\appendix

\chapter{Database Schema}

\section{Lake Table}

\begin{table}[H]
\centering
\caption{Principal lake table fields.}
\label{tab:lake_schema}
\begin{tabularx}{\textwidth}{l l X}
\toprule
Column & Type & Description \\
\midrule
lake\_id & Integer & Unique lake identifier \\
lake\_name & Text & Lake name \\
country & Text & Country \\
lake\_area & Numeric & Lake area \\
elevation & Numeric & Elevation \\
pour\_lat & Numeric & Pour-point latitude \\
pour\_long & Numeric & Pour-point longitude \\
geometry & PostGIS Geometry & Lake boundary \\
lake\_type & Text & Lake classification \\
depth\_avg & Numeric & Average depth \\
vol\_total & Numeric & Total volume \\
wshd\_area & Numeric & Watershed area \\
\bottomrule
\end{tabularx}
\end{table}

\section{Lake Images}

\begin{table}[H]
\centering
\caption{Lake imagery metadata.}
\begin{tabularx}{\textwidth}{l l X}
\toprule
Field & Type & Description \\
\midrule
lake\_id & Integer & Source lake \\
year & Integer & Observation/composite year \\
image\_path & Text & Local GeoTIFF path \\
status & Text & Processing status \\
collection\_used & Text & Sentinel collection used \\
\bottomrule
\end{tabularx}
\end{table}

\section{Lake Tiles}

\begin{table}[H]
\centering
\caption{Intermediate imagery-tile metadata.}
\begin{tabularx}{\textwidth}{l l X}
\toprule
Field & Type & Description \\
\midrule
tile\_id & Integer & Tile identifier \\
lake\_id & Integer & Source lake \\
year & Integer & Composite year \\
tile\_index & Text & Grid position/index \\
tile\_geometry & Geometry & Tile polygon \\
bbox & Geometry/Text & Export bounding box \\
status & Text & Tile state \\
image\_path & Text & Local tile path \\
created\_at & Timestamp & Creation time \\
\bottomrule
\end{tabularx}
\end{table}

\section{Sub-Regions}

\begin{table}[H]
\centering
\caption{Analysis-cell representation.}
\begin{tabularx}{\textwidth}{l l X}
\toprule
Field & Type & Description \\
\midrule
sub\_region\_id & Integer & Unique cell identifier \\
lake\_id & Integer & Parent lake \\
year & Integer & Associated imagery year \\
geometry & Geometry & Analysis-cell polygon \\
coverage\_ratio & Numeric & Lake coverage within cell \\
\bottomrule
\end{tabularx}
\end{table}

\section{Embeddings}

Each valid analysis cell has an associated Prithvi embedding with
768 numerical dimensions. The vector is stored with sufficient
metadata to recover its lake, year, and spatial region.

% ============================================================
% APPENDIX B
% ============================================================

\chapter{Data Processing State Machine}

The processing state of an individual lake-year-tile follows the
general pattern:

\[
Pending
\rightarrow
Processing
\rightarrow
Completed
\]

or:

\[
Processing
\rightarrow
Failed
\]

A previously completed item can be recognized and skipped:

\[
Completed + Existing\ Output
\rightarrow
Skipped
\]

This stateful design is important for large Earth Engine jobs because
a failure in one lake-year should not invalidate successfully
completed work.

% ============================================================
% APPENDIX C
% ============================================================

\chapter{Representative API Structure}

\begin{table}[H]
\centering
\caption{Representative EcoTwin API groups.}
\label{tab:api}
\begin{tabularx}{\textwidth}{l l X}
\toprule
Method & Endpoint family & Purpose \\
\midrule
GET & Lake lookup & Retrieve lake metadata and geometry \\
GET & Similarity query & Retrieve analogous ecosystems \\
GET & Temporal profile & Retrieve NDVI/NDWI/NBR history \\
GET & Forecast & Retrieve analog-based forecast information \\
GET & Report & Retrieve report-generation information \\
GET & Health/status & Check service availability \\
\bottomrule
\end{tabularx}
\end{table}

% ============================================================
% APPENDIX D
% ============================================================

\chapter{Database Migration Procedure}

The project requires safe consolidation of data generated by
different team members.

\section{Backup}

Each developer creates a PostgreSQL custom-format dump:

\begin{lstlisting}
pg_dump \
  -h localhost \
  -U postgres \
  -d ecotwin \
  -Fc \
  -f member.dump
\end{lstlisting}

The custom dump preserves PostgreSQL-native structures including
PostGIS geometries and pgvector embeddings.

\section{Temporary Restore}

A temporary database is created:

\begin{lstlisting}
createdb -h localhost -U postgres ecotwin_member
\end{lstlisting}

The dump is restored:

\begin{lstlisting}
pg_restore \
  -h localhost \
  -U postgres \
  -d ecotwin_member \
  member.dump
\end{lstlisting}

\section{ETL Validation}

The ETL process reads from the temporary database and validates:

\begin{itemize}
    \item missing embeddings;
    \item invalid geometries;
    \item duplicate records;
    \item vector dimensions;
    \item lake/year uniqueness;
    \item image-path existence; and
    \item expected row counts.
\end{itemize}

Validated records are then inserted into the production database.

\section{Production Merge}

Typical tables involved in consolidation include:

\begin{itemize}
    \item \code{lake\_embeddings};
    \item \code{lake\_tiles}; and
    \item \code{lake\_images}.
\end{itemize}

Conflict handling can use database uniqueness constraints and
\code{INSERT ... ON CONFLICT} logic where appropriate.

\section{Cleanup}

After validation, the temporary database can be removed:

\begin{lstlisting}
dropdb -h localhost -U postgres ecotwin_member
\end{lstlisting}

% ============================================================
% APPENDIX E
% ============================================================

\chapter{Recommended Final Report Figures}

The following figures should be inserted before final submission:

\begin{enumerate}
    \item EcoTwin system architecture.
    \item HydroLAKES study-area map.
    \item Lake polygon and 1200 m buffer.
    \item Approximately 8.5 km imagery tile grid.
    \item Example Sentinel-2 annual composite.
    \item Merged lake-year GeoTIFF.
    \item 1 km analysis-cell grid and 15\% coverage filtering.
    \item Prithvi embedding pipeline.
    \item PostgreSQL/PostGIS/pgvector schema.
    \item Similarity-search workflow.
    \item Query lake and Top-$K$ analog map.
    \item Analog lake profile cards.
    \item NDVI temporal comparison.
    \item NDWI temporal comparison.
    \item NBR temporal comparison.
    \item Analog-based forecast plot.
    \item Frontend screenshots.
    \item PDF/report-generation screenshot.
\end{enumerate}

\chapter{Final Experimental Tables}

Before submission, populate the following values directly from the
final production database and pipeline logs:

\begin{table}[H]
\centering
\caption{Final production statistics to be reported.}
\begin{tabular}{lr}
\toprule
Metric & Final value \\
\midrule
Source lakes & TBD \\
Processed lakes & TBD \\
Lake-years attempted & TBD \\
Lake-years completed & TBD \\
Imagery tiles generated & TBD \\
Merged lake images & TBD \\
Analysis cells generated & TBD \\
Valid analysis cells & TBD \\
Embeddings generated & TBD \\
Embedding failures & TBD \\
Indexed vectors & TBD \\
Temporal records & TBD \\
\bottomrule
\end{tabular}
\end{table}

\begin{table}[H]
\centering
\caption{Final performance metrics to be reported.}
\begin{tabular}{lr}
\toprule
Operation & Final measurement \\
\midrule
Average imagery extraction time & TBD \\
Average raster merge time & TBD \\
Average embedding time/cell & TBD \\
Similarity query latency & TBD \\
Temporal retrieval latency & TBD \\
Report generation time & TBD \\
Total preprocessing time & TBD \\
\bottomrule
\end{tabular}
\end{table}

% ============================================================
% END DOCUMENT
% ============================================================

\end{document}